% Generated draft: requires mathematical and source review.
\documentclass[11pt]{article}
\usepackage[margin=25mm]{geometry}
\usepackage{fontspec}
\setmainfont{Latin Modern Roman}
\usepackage{amsmath,amssymb,mathtools,unicode-math}
\setmathfont{Latin Modern Math}
\usepackage{enumitem,needspace,xurl,hyperref}
\setlength{\parindent}{0pt}
\setlength{\parskip}{5pt}
\setlength{\emergencystretch}{3em}
\newcommand{\N}{\mathbb{N}}
\newcommand{\Z}{\mathbb{Z}}
\newcommand{\Q}{\mathbb{Q}}
\newcommand{\R}{\mathbb{R}}
\newcommand{\C}{\mathbb{C}}
\newcommand{\F}{\mathbb{F}}
\newcommand{\A}{\mathbb{A}}
\newcommand{\PP}{\mathbb{P}}
\newcommand{\E}{\mathbb{E}}
\newcommand{\eps}{\varepsilon}
\newcommand{\dd}{\,\mathrm{d}}
\begin{document}
\section*{UnsolvedMath ID: 30005649}
\textbf{Source problem code:} catalogue entry
\subsection{Definitions and mathematical statement}
Let \(p\) be a prime and let \(n\geq 3\) be an integer. The supplied statement denotes a ring of Witt vectors by \(W\), but does not specify its coefficient field or the hypotheses on that field. Write \(k\) for the intended residue field and \(\overline G=G\times_{\operatorname{Spec}W}\operatorname{Spec}k\) for the special fiber. The choice of \(W\) and \(k\) remains an unresolved part of the formulation.

Let \(G\) be a finite flat commutative group scheme over \(W\), of constant rank \(p^{2n}\), annihilated by \(p\), so that \([p]_G=0\). Assume that \(\overline G\) is quasi-supersingular in the sense intended by the source. The supplied record does not define quasi-supersingularity, so this hypothesis cannot be made explicit from the available evidence.

For a finite locally free commutative group scheme \(H\) over a base scheme \(S\), its Cartier dual \(H^D\) represents the sheaf \(\underline{\operatorname{Hom}}_S(H,\mathbb G_{m,S})\). Interpret self-duality as the existence of an isomorphism \(H\cong H^D\) over \(S\), with no additional compatibility stipulated.

\textbf{Provisional question from the supplied clean statement.} For every such \(G\), must there exist isomorphisms
\[
G\cong G^D\quad\text{over }W,
\qquad
\overline G\cong(\overline G)^D\quad\text{over }k?
\]
The intended coefficient field, the meaning of quasi-supersingularity, and the precise meaning of self-duality require confirmation. Moreover, the supplied original excerpt discusses equivalence of unspecified conditions (1) and (2) in Proposition 1 under an additional self-duality assumption; it does not itself state the universal question above. Thus the correspondence between this question and the original report remains unresolved.
\subsection{Short English statement}
The supplied clean statement asks whether every finite flat commutative group scheme of rank \(p^{2n}\), killed by \(p\), over the intended Witt-vector ring, with quasi-supersingular special fiber, is self-dual together with its special fiber when \(n\geq3\). The base field, quasi-supersingularity hypothesis, and correspondence with the supplied original excerpt need clarification.
\subsection{Sources}
\begin{itemize}
\item Local UnsolvedMath registry, permanent ID \(30005649\), exact source problem code \texttt{catalogue entry}, title \emph{Self-Duality of Quasi-Supersingular Group Schemes}. The supplied record labels the statement \texttt{corrected\_verified} and reports comparison with the official report; this is supplied provenance, not independent verification here. No individual proposer is supplied.
\item Supplied report citation: \emph{MFO--RIMS Tandem Workshop: Arithmetic Homotopy and Galois Theory (2024)}, Oberwolfach Reports, DOI 10.4171/owr/2023/42, pp.\ 2377--2488. \url{https://doi.org/10.4171/owr/2023/42}. The record separately gives proposed year 2023 and an \texttt{owr-year:2023} tag; the relationship to the citation's 2024 date is unspecified.
\item Supplied literature source: ``Takao, Coleman conjecture report.'' \url{https://doi.org/10.4171/owr/2023/42}.
\item Supplied literature source: ``Hoshi, finite-flat group-scheme results.'' \url{https://www.kurims.kyoto-u.ac.jp/~yuichiro/papers_e.html}. No specific paper or theorem is identified in the supplied record.
\item Literature assessment checked at \texttt{2026-08-22}: supplied dated curation, not independently verified. The record labels the problem partially solved and states that the genus-two case of an associated supersingular torsion-point problem is proved affirmative, while higher-genus Raynaud-type cases remain open strategies and questions in the report. This assessment concerns an associated problem; its precise implications for the self-duality question are not supplied. The record describes this assessment as dated literature triage derived from the supplied OWR curation and requests independent verification of time-sensitive status and cited claims.
\end{itemize}
\end{document}
